% --- Archon physics formalization source begin ---
% archon:physics
% archon:covers PhyXMiniProblems/problem_phyx_mini_0699.lean
% archon:source-report reports/phyx_mini/problem_phyx_mini_0699.source.json

\chapter{Physics problem phyx\_mini\_0699}
\label{ch:PhyXMiniProblems_problem_phyx_mini_0699}

\paragraph{Problem source.}
\#\# Physical scenario

A \textbackslash{}(1.0\textbackslash{},\textbackslash{}mathrm\{m\}\textbackslash{})-long, \textbackslash{}(200\textbackslash{},\textbackslash{}mathrm\{g\}\textbackslash{}) rod is hinged at one end and connected to a wall. It is held out horizontally, then released.

\#\# Question

What is the speed of the tip of the rod as it hits the wall?

\#\# Answer choices

- A: A: 5.2m/s

- B: B: 5.6m/s

- C: C: 5.4m/s

- D: D: 5.8m/s

\#\# Figure caption (auxiliary; use the image as primary evidence)

The image shows a diagram of a thin rod hinged at one end, forming an L shape with the x and y axes. Key elements include:

- **Rod and Hinge**: 
  - The rod is hinged at the origin where the x and y axes intersect. It initially stands vertically along the y-axis and rotates downward, sweeping towards the x-axis.
  - The rod has circular markings indicating its rotation.

- **Dimensions and Properties**:
  - The length of the rod, \textbackslash{}( L \textbackslash{}), is 1.0 m.
  - Initial conditions are \textbackslash{}( y\_\{cm0\} = 0 \textbackslash{}), \textbackslash{}(\textbackslash{}omega\_0 = 0 \textbackslash{}, \textbackslash{}text\{rad/s\}\textbackslash{}), and mass \textbackslash{}( m = 0.20 \textbackslash{}, \textbackslash{}text\{kg\}\textbackslash{}).

- **After Rotation**:
  - The center of mass after rotation is at \textbackslash{}( y\_\{cm1\} = -\textbackslash{}frac\{1\}\{2\}L \textbackslash{}).

- **Calculations**:
  - The goal is to find the tip velocity, \textbackslash{}( v\_\{tip\} \textbackslash{}), using the relationship \textbackslash{}( v\_\{tip\} = \textbackslash{}omega\_1 L \textbackslash{}).

- **Arrows**:
  - A green arrow labeled \textbackslash{}( \textbackslash{}vec\{v\}\_\{tip\} \textbackslash{}) indicates the velocity direction at the tip of the rod.
  - A grey arrow shows the rotational motion direction of the rod.

This diagram is likely used in a physics context to explain rotational motion and center of mass displacement.

\#\# Dataset metadata

Rotational Motion; Physical Model Grounding Reasoning, Multi-Formula Reasoning

\paragraph{Recorded answer/context.}
C: C: 5.4m/s

\paragraph{Figure/image path.}
/root/proposal\_for\_physic/science-mango/phyx\_mini\_run/phyx\_data/test\_image/699.png

\paragraph{Formalization target.}
create a compiling Lean file with sorry bodies at `PhyXMiniProblems/problem\_phyx\_mini\_0699.lean`. The Lean declarations must preserve the physical quantities, dimensions or dimensional roles, figure labels, governing-law hypotheses, and final relation expressed by this problem.
Use Mathlib/Physlib names found through LeanExplore where available. If a physics API is missing, introduce faithful local abstractions rather than scalar placeholder aliases.

\begin{theorem}[Physics formalization target]
\label{thm:physics:phyx_mini_0699:target}
\lean{PhyXMiniProblems.ProblemPhyXMini0699.problem_phyx_mini_0699}
\uses{def:physics:phyx-mini-0699:phyxminiproblems-problemphyxmini0699-speedinmeterspersecond, def:physics:phyx-mini-0699:phyxminiproblems-problemphyxmini0699-hingedrodsetup, def:physics:phyx-mini-0699:phyxminiproblems-problemphyxmini0699-matchesprimaryhingedrodfigure, def:physics:phyx-mini-0699:phyxminiproblems-problemphyxmini0699-matcheshingedrodproblemdata, def:physics:phyx-mini-0699:phyxminiproblems-problemphyxmini0699-matcheshingedrodgeometry, def:physics:phyx-mini-0699:phyxminiproblems-problemphyxmini0699-hasphysicalhingedrodparameters, def:physics:phyx-mini-0699:phyxminiproblems-problemphyxmini0699-satisfiesfallinghingedrodlaws, def:physics:phyx-mini-0699:phyxminiproblems-problemphyxmini0699-answerchoice, def:physics:phyx-mini-0699:phyxminiproblems-problemphyxmini0699-recordedanswerchoice, def:physics:phyx-mini-0699:phyxminiproblems-problemphyxmini0699-matchesdisplayedspeed}
The assigned autoformalize agent should translate this physics problem into Lean declarations in the covered file, with theorem and lemma proof bodies written as `by sorry`.
\end{theorem}
\begin{proof}
This is an autoformalization task, not a proof task. Produce statements that can later be proved by the physics prover without weakening the physical model.
\end{proof}

% --- Archon declaration topology begin ---
\section{Declaration topology}
\begin{definition}[acceleration Dimension]
  \label{def:physics:phyx-mini-0699:phyxminiproblems-problemphyxmini0699-accelerationdimension}
  \lean{PhyXMiniProblems.ProblemPhyXMini0699.accelerationDimension}
  The physical dimension of acceleration, L T⁻²
\end{definition}
\begin{proof}
  This is the declared model definition of acceleration Dimension.
\end{proof}
\begin{definition}[Length Quantity]
  \label{def:physics:phyx-mini-0699:phyxminiproblems-problemphyxmini0699-lengthquantity}
  \lean{PhyXMiniProblems.ProblemPhyXMini0699.LengthQuantity}
  A nonnegative physical length magnitude
\end{definition}
\begin{proof}
  This is the declared model definition of Length Quantity.
\end{proof}
\begin{definition}[Signed Length Quantity]
  \label{def:physics:phyx-mini-0699:phyxminiproblems-problemphyxmini0699-signedlengthquantity}
  \lean{PhyXMiniProblems.ProblemPhyXMini0699.SignedLengthQuantity}
  A signed physical length, used for Cartesian coordinates and heights
\end{definition}
\begin{proof}
  This is the declared model definition of Signed Length Quantity.
\end{proof}
\begin{definition}[Mass Quantity]
  \label{def:physics:phyx-mini-0699:phyxminiproblems-problemphyxmini0699-massquantity}
  \lean{PhyXMiniProblems.ProblemPhyXMini0699.MassQuantity}
  A nonnegative physical mass
\end{definition}
\begin{proof}
  This is the declared model definition of Mass Quantity.
\end{proof}
\begin{definition}[Acceleration Quantity]
  \label{def:physics:phyx-mini-0699:phyxminiproblems-problemphyxmini0699-accelerationquantity}
  \lean{PhyXMiniProblems.ProblemPhyXMini0699.AccelerationQuantity}
  \uses{def:physics:phyx-mini-0699:phyxminiproblems-problemphyxmini0699-accelerationdimension}
  A nonnegative gravitational-acceleration magnitude
\end{definition}
\begin{proof}
  This is the declared model definition of Acceleration Quantity.
\end{proof}
\begin{definition}[Angular Speed Quantity]
  \label{def:physics:phyx-mini-0699:phyxminiproblems-problemphyxmini0699-angularspeedquantity}
  \lean{PhyXMiniProblems.ProblemPhyXMini0699.AngularSpeedQuantity}
  A nonnegative angular-speed magnitude, with dimension T⁻¹
\end{definition}
\begin{proof}
  This is the declared model definition of Angular Speed Quantity.
\end{proof}
\begin{definition}[Moment Of Inertia Quantity]
  \label{def:physics:phyx-mini-0699:phyxminiproblems-problemphyxmini0699-momentofinertiaquantity}
  \lean{PhyXMiniProblems.ProblemPhyXMini0699.MomentOfInertiaQuantity}
  A nonnegative moment of inertia about an axis, with dimension M L²
\end{definition}
\begin{proof}
  This is the declared model definition of Moment Of Inertia Quantity.
\end{proof}
\begin{definition}[Speed Quantity]
  \label{def:physics:phyx-mini-0699:phyxminiproblems-problemphyxmini0699-speedquantity}
  \lean{PhyXMiniProblems.ProblemPhyXMini0699.SpeedQuantity}
  A nonnegative physical speed magnitude
\end{definition}
\begin{proof}
  This is the declared model definition of Speed Quantity.
\end{proof}
\begin{definition}[Energy Quantity]
  \label{def:physics:phyx-mini-0699:phyxminiproblems-problemphyxmini0699-energyquantity}
  \lean{PhyXMiniProblems.ProblemPhyXMini0699.EnergyQuantity}
  A physical energy, with dimension M L² T⁻²
\end{definition}
\begin{proof}
  This is the declared model definition of Energy Quantity.
\end{proof}
\begin{definition}[nonnegative SIReadout]
  \label{def:physics:phyx-mini-0699:phyxminiproblems-problemphyxmini0699-nonnegativesireadout}
  \lean{PhyXMiniProblems.ProblemPhyXMini0699.nonnegativeSIReadout}
  Read a nonnegative dimensionful quantity in coherent SI units
\end{definition}
\begin{proof}
  This is the declared model definition of nonnegative SIReadout.
\end{proof}
\begin{definition}[length In Meters]
  \label{def:physics:phyx-mini-0699:phyxminiproblems-problemphyxmini0699-lengthinmeters}
  \lean{PhyXMiniProblems.ProblemPhyXMini0699.lengthInMeters}
  \uses{def:physics:phyx-mini-0699:phyxminiproblems-problemphyxmini0699-lengthquantity, def:physics:phyx-mini-0699:phyxminiproblems-problemphyxmini0699-nonnegativesireadout}
  Metre readout of a nonnegative physical length
\end{definition}
\begin{proof}
  This is the declared model definition of length In Meters.
\end{proof}
\begin{definition}[signed Length In Meters]
  \label{def:physics:phyx-mini-0699:phyxminiproblems-problemphyxmini0699-signedlengthinmeters}
  \lean{PhyXMiniProblems.ProblemPhyXMini0699.signedLengthInMeters}
  \uses{def:physics:phyx-mini-0699:phyxminiproblems-problemphyxmini0699-signedlengthquantity}
  Metre readout of a signed physical coordinate or height
\end{definition}
\begin{proof}
  This is the declared model definition of signed Length In Meters.
\end{proof}
\begin{definition}[mass In Kilograms]
  \label{def:physics:phyx-mini-0699:phyxminiproblems-problemphyxmini0699-massinkilograms}
  \lean{PhyXMiniProblems.ProblemPhyXMini0699.massInKilograms}
  \uses{def:physics:phyx-mini-0699:phyxminiproblems-problemphyxmini0699-massquantity, def:physics:phyx-mini-0699:phyxminiproblems-problemphyxmini0699-nonnegativesireadout}
  Kilogram readout of a physical mass
\end{definition}
\begin{proof}
  This is the declared model definition of mass In Kilograms.
\end{proof}
\begin{definition}[acceleration In Meters Per Second Squared]
  \label{def:physics:phyx-mini-0699:phyxminiproblems-problemphyxmini0699-accelerationinmeterspersecondsquared}
  \lean{PhyXMiniProblems.ProblemPhyXMini0699.accelerationInMetersPerSecondSquared}
  \uses{def:physics:phyx-mini-0699:phyxminiproblems-problemphyxmini0699-accelerationquantity, def:physics:phyx-mini-0699:phyxminiproblems-problemphyxmini0699-nonnegativesireadout}
  Metre-per-second-squared readout of an acceleration magnitude
\end{definition}
\begin{proof}
  This is the declared model definition of acceleration In Meters Per Second Squared.
\end{proof}
\begin{definition}[angular Speed In Radians Per Second]
  \label{def:physics:phyx-mini-0699:phyxminiproblems-problemphyxmini0699-angularspeedinradianspersecond}
  \lean{PhyXMiniProblems.ProblemPhyXMini0699.angularSpeedInRadiansPerSecond}
  \uses{def:physics:phyx-mini-0699:phyxminiproblems-problemphyxmini0699-angularspeedquantity, def:physics:phyx-mini-0699:phyxminiproblems-problemphyxmini0699-nonnegativesireadout}
  Radian-per-second readout of an angular-speed magnitude
\end{definition}
\begin{proof}
  This is the declared model definition of angular Speed In Radians Per Second.
\end{proof}
\begin{definition}[moment Of Inertia In Kilogram Meters Squared]
  \label{def:physics:phyx-mini-0699:phyxminiproblems-problemphyxmini0699-momentofinertiainkilogrammeterssquared}
  \lean{PhyXMiniProblems.ProblemPhyXMini0699.momentOfInertiaInKilogramMetersSquared}
  \uses{def:physics:phyx-mini-0699:phyxminiproblems-problemphyxmini0699-momentofinertiaquantity, def:physics:phyx-mini-0699:phyxminiproblems-problemphyxmini0699-nonnegativesireadout}
  Kilogram-metre-squared readout of a moment of inertia
\end{definition}
\begin{proof}
  This is the declared model definition of moment Of Inertia In Kilogram Meters Squared.
\end{proof}
\begin{definition}[speed In Meters Per Second]
  \label{def:physics:phyx-mini-0699:phyxminiproblems-problemphyxmini0699-speedinmeterspersecond}
  \lean{PhyXMiniProblems.ProblemPhyXMini0699.speedInMetersPerSecond}
  \uses{def:physics:phyx-mini-0699:phyxminiproblems-problemphyxmini0699-speedquantity}
  Metre-per-second readout of a physical speed magnitude
\end{definition}
\begin{proof}
  This is the declared model definition of speed In Meters Per Second.
\end{proof}
\begin{definition}[energy In Joules]
  \label{def:physics:phyx-mini-0699:phyxminiproblems-problemphyxmini0699-energyinjoules}
  \lean{PhyXMiniProblems.ProblemPhyXMini0699.energyInJoules}
  \uses{def:physics:phyx-mini-0699:phyxminiproblems-problemphyxmini0699-energyquantity}
  Joule readout of a physical energy
\end{definition}
\begin{proof}
  This is the declared model definition of energy In Joules.
\end{proof}
\begin{definition}[Motion Event]
  \label{def:physics:phyx-mini-0699:phyxminiproblems-problemphyxmini0699-motionevent}
  \lean{PhyXMiniProblems.ProblemPhyXMini0699.MotionEvent}
  The two rod configurations compared by conservation of energy
\end{definition}
\begin{proof}
  This is the declared model definition of Motion Event.
\end{proof}
\begin{definition}[Cartesian Axis]
  \label{def:physics:phyx-mini-0699:phyxminiproblems-problemphyxmini0699-cartesianaxis}
  \lean{PhyXMiniProblems.ProblemPhyXMini0699.CartesianAxis}
  The axes explicitly labelled in the supplied image
\end{definition}
\begin{proof}
  This is the declared model definition of Cartesian Axis.
\end{proof}
\begin{definition}[Rod Point]
  \label{def:physics:phyx-mini-0699:phyxminiproblems-problemphyxmini0699-rodpoint}
  \lean{PhyXMiniProblems.ProblemPhyXMini0699.RodPoint}
  Distinguished points of the idealized rod
\end{definition}
\begin{proof}
  This is the declared model definition of Rod Point.
\end{proof}
\begin{definition}[Rod Pose]
  \label{def:physics:phyx-mini-0699:phyxminiproblems-problemphyxmini0699-rodpose}
  \lean{PhyXMiniProblems.ProblemPhyXMini0699.RodPose}
  The two physical orientations shown in the before/after diagram
\end{definition}
\begin{proof}
  This is the declared model definition of Rod Pose.
\end{proof}
\begin{definition}[Figure Label]
  \label{def:physics:phyx-mini-0699:phyxminiproblems-problemphyxmini0699-figurelabel}
  \lean{PhyXMiniProblems.ProblemPhyXMini0699.FigureLabel}
  Text and mathematical labels visibly printed in the primary image
\end{definition}
\begin{proof}
  This is the declared model definition of Figure Label.
\end{proof}
\begin{definition}[Rotation Sense]
  \label{def:physics:phyx-mini-0699:phyxminiproblems-problemphyxmini0699-rotationsense}
  \lean{PhyXMiniProblems.ProblemPhyXMini0699.RotationSense}
  Sense of the grey quarter-turn arrow in the image
\end{definition}
\begin{proof}
  This is the declared model definition of Rotation Sense.
\end{proof}
\begin{definition}[Tip Velocity Direction]
  \label{def:physics:phyx-mini-0699:phyxminiproblems-problemphyxmini0699-tipvelocitydirection}
  \lean{PhyXMiniProblems.ProblemPhyXMini0699.TipVelocityDirection}
  Direction of the green tip-velocity arrow at wall impact
\end{definition}
\begin{proof}
  This is the declared model definition of Tip Velocity Direction.
\end{proof}
\begin{definition}[Rod Mass Distribution]
  \label{def:physics:phyx-mini-0699:phyxminiproblems-problemphyxmini0699-rodmassdistribution}
  \lean{PhyXMiniProblems.ProblemPhyXMini0699.RodMassDistribution}
  Idealized mass distribution of the pictured rod
\end{definition}
\begin{proof}
  This is the declared model definition of Rod Mass Distribution.
\end{proof}
\begin{definition}[Hinge Model]
  \label{def:physics:phyx-mini-0699:phyxminiproblems-problemphyxmini0699-hingemodel}
  \lean{PhyXMiniProblems.ProblemPhyXMini0699.HingeModel}
  Constraint model at the wall attachment
\end{definition}
\begin{proof}
  This is the declared model definition of Hinge Model.
\end{proof}
\begin{definition}[Gravity Model]
  \label{def:physics:phyx-mini-0699:phyxminiproblems-problemphyxmini0699-gravitymodel}
  \lean{PhyXMiniProblems.ProblemPhyXMini0699.GravityModel}
  Gravity model used during the fall
\end{definition}
\begin{proof}
  This is the declared model definition of Gravity Model.
\end{proof}
\begin{definition}[Hinged Rod Figure]
  \label{def:physics:phyx-mini-0699:phyxminiproblems-problemphyxmini0699-hingedrodfigure}
  \lean{PhyXMiniProblems.ProblemPhyXMini0699.HingedRodFigure}
  \uses{def:physics:phyx-mini-0699:phyxminiproblems-problemphyxmini0699-motionevent, def:physics:phyx-mini-0699:phyxminiproblems-problemphyxmini0699-cartesianaxis, def:physics:phyx-mini-0699:phyxminiproblems-problemphyxmini0699-rodpose, def:physics:phyx-mini-0699:phyxminiproblems-problemphyxmini0699-figurelabel, def:physics:phyx-mini-0699:phyxminiproblems-problemphyxmini0699-rotationsense, def:physics:phyx-mini-0699:phyxminiproblems-problemphyxmini0699-tipvelocitydirection}
  Qualitative and label evidence read directly from the supplied bitmap
\end{definition}
\begin{proof}
  This is the declared model definition of Hinged Rod Figure.
\end{proof}
\begin{definition}[Hinged Rod Setup]
  \label{def:physics:phyx-mini-0699:phyxminiproblems-problemphyxmini0699-hingedrodsetup}
  \lean{PhyXMiniProblems.ProblemPhyXMini0699.HingedRodSetup}
  \uses{def:physics:phyx-mini-0699:phyxminiproblems-problemphyxmini0699-lengthquantity, def:physics:phyx-mini-0699:phyxminiproblems-problemphyxmini0699-signedlengthquantity, def:physics:phyx-mini-0699:phyxminiproblems-problemphyxmini0699-massquantity, def:physics:phyx-mini-0699:phyxminiproblems-problemphyxmini0699-accelerationquantity, def:physics:phyx-mini-0699:phyxminiproblems-problemphyxmini0699-angularspeedquantity, def:physics:phyx-mini-0699:phyxminiproblems-problemphyxmini0699-momentofinertiaquantity, def:physics:phyx-mini-0699:phyxminiproblems-problemphyxmini0699-speedquantity, def:physics:phyx-mini-0699:phyxminiproblems-problemphyxmini0699-energyquantity, def:physics:phyx-mini-0699:phyxminiproblems-problemphyxmini0699-motionevent, def:physics:phyx-mini-0699:phyxminiproblems-problemphyxmini0699-cartesianaxis, def:physics:phyx-mini-0699:phyxminiproblems-problemphyxmini0699-rodpoint, def:physics:phyx-mini-0699:phyxminiproblems-problemphyxmini0699-rodpose, def:physics:phyx-mini-0699:phyxminiproblems-problemphyxmini0699-rodmassdistribution, def:physics:phyx-mini-0699:phyxminiproblems-problemphyxmini0699-hingemodel, def:physics:phyx-mini-0699:phyxminiproblems-problemphyxmini0699-gravitymodel, def:physics:phyx-mini-0699:phyxminiproblems-problemphyxmini0699-hingedrodfigure}
  Independent physical quantities and geometry of the experiment. Point coordinates are signed dimensionful lengths relative to the hinge. The inertia, energies, angular speeds, and impact-tip speed are stored independently; their relations are supplied only by the governing laws
\end{definition}
\begin{proof}
  This is the declared model definition of Hinged Rod Setup.
\end{proof}
\begin{definition}[Matches Primary Hinged Rod Figure]
  \label{def:physics:phyx-mini-0699:phyxminiproblems-problemphyxmini0699-matchesprimaryhingedrodfigure}
  \lean{PhyXMiniProblems.ProblemPhyXMini0699.MatchesPrimaryHingedRodFigure}
  \uses{def:physics:phyx-mini-0699:phyxminiproblems-problemphyxmini0699-hingedrodsetup}
  Primary-image evidence. The bitmap, unlike one sentence of the auxiliary caption, clearly shows the release rod along +x and the impact rod along -y. It also shows the hinge at the axes' intersection, center-of-mass markers in both poses, clockwise rotation, and a leftward impact velocity
\end{definition}
\begin{proof}
  This is the declared model definition of Matches Primary Hinged Rod Figure.
\end{proof}
\begin{definition}[Matches Hinged Rod Problem Data]
  \label{def:physics:phyx-mini-0699:phyxminiproblems-problemphyxmini0699-matcheshingedrodproblemdata}
  \lean{PhyXMiniProblems.ProblemPhyXMini0699.MatchesHingedRodProblemData}
  \uses{def:physics:phyx-mini-0699:phyxminiproblems-problemphyxmini0699-lengthinmeters, def:physics:phyx-mini-0699:phyxminiproblems-problemphyxmini0699-signedlengthinmeters, def:physics:phyx-mini-0699:phyxminiproblems-problemphyxmini0699-massinkilograms, def:physics:phyx-mini-0699:phyxminiproblems-problemphyxmini0699-accelerationinmeterspersecondsquared, def:physics:phyx-mini-0699:phyxminiproblems-problemphyxmini0699-angularspeedinradianspersecond, def:physics:phyx-mini-0699:phyxminiproblems-problemphyxmini0699-hingedrodsetup}
  Numerical and qualitative problem data. The diagram gives L = 1.0 m, m = 0.20 kg, y\_cm,0 = 0, and ω₀ = 0; standard near-Earth gravity is calibrated as 9.8 m/s². No impact angular speed or tip speed occurs here
\end{definition}
\begin{proof}
  This is the declared model definition of Matches Hinged Rod Problem Data.
\end{proof}
\begin{definition}[Matches Hinged Rod Geometry]
  \label{def:physics:phyx-mini-0699:phyxminiproblems-problemphyxmini0699-matcheshingedrodgeometry}
  \lean{PhyXMiniProblems.ProblemPhyXMini0699.MatchesHingedRodGeometry}
  \uses{def:physics:phyx-mini-0699:phyxminiproblems-problemphyxmini0699-lengthinmeters, def:physics:phyx-mini-0699:phyxminiproblems-problemphyxmini0699-signedlengthinmeters, def:physics:phyx-mini-0699:phyxminiproblems-problemphyxmini0699-hingedrodsetup}
  Physical coordinates implied by a length-L uniform rod hinged at the origin. At release the tip is (L, 0) and the center of mass is (L/2, 0); at impact they are (0, -L) and (0, -L/2) respectively. These are figure and uniform-rod geometry, not the requested speed
\end{definition}
\begin{proof}
  This is the declared model definition of Matches Hinged Rod Geometry.
\end{proof}
\begin{definition}[Has Physical Hinged Rod Parameters]
  \label{def:physics:phyx-mini-0699:phyxminiproblems-problemphyxmini0699-hasphysicalhingedrodparameters}
  \lean{PhyXMiniProblems.ProblemPhyXMini0699.HasPhysicalHingedRodParameters}
  \uses{def:physics:phyx-mini-0699:phyxminiproblems-problemphyxmini0699-lengthinmeters, def:physics:phyx-mini-0699:phyxminiproblems-problemphyxmini0699-massinkilograms, def:physics:phyx-mini-0699:phyxminiproblems-problemphyxmini0699-accelerationinmeterspersecondsquared, def:physics:phyx-mini-0699:phyxminiproblems-problemphyxmini0699-momentofinertiainkilogrammeterssquared, def:physics:phyx-mini-0699:phyxminiproblems-problemphyxmini0699-hingedrodsetup}
  Positivity and nondegeneracy of the physical rod and gravity
\end{definition}
\begin{proof}
  This is the declared model definition of Has Physical Hinged Rod Parameters.
\end{proof}
\begin{definition}[Satisfies Falling Hinged Rod Laws]
  \label{def:physics:phyx-mini-0699:phyxminiproblems-problemphyxmini0699-satisfiesfallinghingedrodlaws}
  \lean{PhyXMiniProblems.ProblemPhyXMini0699.SatisfiesFallingHingedRodLaws}
  \uses{def:physics:phyx-mini-0699:phyxminiproblems-problemphyxmini0699-lengthinmeters, def:physics:phyx-mini-0699:phyxminiproblems-problemphyxmini0699-signedlengthinmeters, def:physics:phyx-mini-0699:phyxminiproblems-problemphyxmini0699-massinkilograms, def:physics:phyx-mini-0699:phyxminiproblems-problemphyxmini0699-accelerationinmeterspersecondsquared, def:physics:phyx-mini-0699:phyxminiproblems-problemphyxmini0699-angularspeedinradianspersecond, def:physics:phyx-mini-0699:phyxminiproblems-problemphyxmini0699-momentofinertiainkilogrammeterssquared, def:physics:phyx-mini-0699:phyxminiproblems-problemphyxmini0699-speedinmeterspersecond, def:physics:phyx-mini-0699:phyxminiproblems-problemphyxmini0699-energyinjoules, def:physics:phyx-mini-0699:phyxminiproblems-problemphyxmini0699-hingedrodsetup}
  Governing relations for a uniform rod rotating about a fixed hinge: the hinge-axis moment is I = m L² / 3; gravitational potential energy is m g y\_cm; rotational energy is Physlib's RigidBody.rotationalKineticEnergy for the angular-velocity vector and hinge-referenced inertia tensor; mechanical energy is conserved from release to wall impact; fixed-axis kinematics gives v\_tip = ω₁ L. None of these fields states the square-root speed, 5.4 m/s, or an answer choice
\end{definition}
\begin{proof}
  This is the declared model definition of Satisfies Falling Hinged Rod Laws.
\end{proof}
\begin{lemma}[impact tip speed squared]
  \label{lem:physics:phyx-mini-0699:phyxminiproblems-problemphyxmini0699-impact-tip-speed-squared}
  \lean{PhyXMiniProblems.ProblemPhyXMini0699.impact_tip_speed_squared}
  \uses{def:physics:phyx-mini-0699:phyxminiproblems-problemphyxmini0699-lengthinmeters, def:physics:phyx-mini-0699:phyxminiproblems-problemphyxmini0699-accelerationinmeterspersecondsquared, def:physics:phyx-mini-0699:phyxminiproblems-problemphyxmini0699-speedinmeterspersecond, def:physics:phyx-mini-0699:phyxminiproblems-problemphyxmini0699-hingedrodsetup, def:physics:phyx-mini-0699:phyxminiproblems-problemphyxmini0699-matcheshingedrodproblemdata, def:physics:phyx-mini-0699:phyxminiproblems-problemphyxmini0699-matcheshingedrodgeometry, def:physics:phyx-mini-0699:phyxminiproblems-problemphyxmini0699-hasphysicalhingedrodparameters, def:physics:phyx-mini-0699:phyxminiproblems-problemphyxmini0699-satisfiesfallinghingedrodlaws}
  Energy conservation, the center-of-mass drop L/2, the hinge inertia m L²/3, and v\_tip = ω₁ L imply v\_tip² = 3 g L. This is a derived mechanics result rather than a law premise
\end{lemma}
\begin{proof}
  The conclusion follows from the governing relations and figure readouts named in the statement of impact tip speed squared.
\end{proof}
\begin{definition}[Answer Choice]
  \label{def:physics:phyx-mini-0699:phyxminiproblems-problemphyxmini0699-answerchoice}
  \lean{PhyXMiniProblems.ProblemPhyXMini0699.AnswerChoice}
  Labels of the four numerical answer choices
\end{definition}
\begin{proof}
  This is the declared model definition of Answer Choice.
\end{proof}
\begin{definition}[meters Per Second]
  \label{def:physics:phyx-mini-0699:phyxminiproblems-problemphyxmini0699-answerchoice-meterspersecond}
  \lean{PhyXMiniProblems.ProblemPhyXMini0699.AnswerChoice.metersPerSecond}
  \uses{def:physics:phyx-mini-0699:phyxminiproblems-problemphyxmini0699-answerchoice}
  Speed in metres per second printed beside each answer label
\end{definition}
\begin{proof}
  This is the declared model definition of meters Per Second.
\end{proof}
\begin{definition}[recorded Answer Choice]
  \label{def:physics:phyx-mini-0699:phyxminiproblems-problemphyxmini0699-recordedanswerchoice}
  \lean{PhyXMiniProblems.ProblemPhyXMini0699.recordedAnswerChoice}
  \uses{def:physics:phyx-mini-0699:phyxminiproblems-problemphyxmini0699-answerchoice}
  Dataset metadata recording the supplied answer label
\end{definition}
\begin{proof}
  This is the declared model definition of recorded Answer Choice.
\end{proof}
\begin{definition}[displayed Speed Resolution]
  \label{def:physics:phyx-mini-0699:phyxminiproblems-problemphyxmini0699-displayedspeedresolution}
  \lean{PhyXMiniProblems.ProblemPhyXMini0699.displayedSpeedResolution}
  One unit in the final displayed decimal place
\end{definition}
\begin{proof}
  This is the declared model definition of displayed Speed Resolution.
\end{proof}
\begin{definition}[Matches Displayed Speed]
  \label{def:physics:phyx-mini-0699:phyxminiproblems-problemphyxmini0699-matchesdisplayedspeed}
  \lean{PhyXMiniProblems.ProblemPhyXMini0699.MatchesDisplayedSpeed}
  \uses{def:physics:phyx-mini-0699:phyxminiproblems-problemphyxmini0699-speedquantity, def:physics:phyx-mini-0699:phyxminiproblems-problemphyxmini0699-speedinmeterspersecond, def:physics:phyx-mini-0699:phyxminiproblems-problemphyxmini0699-answerchoice, def:physics:phyx-mini-0699:phyxminiproblems-problemphyxmini0699-displayedspeedresolution}
  Agreement with a speed displayed to the nearest tenth of a metre/second
\end{definition}
\begin{proof}
  This is the declared model definition of Matches Displayed Speed.
\end{proof}
% --- Archon declaration topology end ---
% --- Archon physics formalization source end ---
